\section{Conclusion}
\label{sec:conclusion}

We began with LoRA installed correctly on Mamba-130m --- right target modules, right parameter counts, 1.14\% trainable weights --- and watched SST-2 accuracy sit at 50.9\% for three consecutive training epochs, indistinguishable from the zero-shot baseline of 49.1\%. We have now traced what that silence means.

The finding is not that LoRA fails in some general sense. The finding is that projection-only LoRA fails on Mamba for a specific and mechanistically coherent reason: the selective state space scan, implemented as a custom CUDA parallel associative scan, sits between the classification head and every LoRA weight matrix in every layer. For next-token prediction --- the task this kernel was built for --- gradient flows through it naturally. For a randomly-initialized classification head attached to the final hidden state, the gradient exists numerically but carries no task-discriminative information to the projection layers. The loss oscillates. Accuracy does not move. And on MNLI, it moves in the wrong direction.

This is the distinction the field needs: \textbf{API compatibility is not gradient compatibility}. Mamba's projection layers are \texttt{nn.Linear}. LoRA installs on them without complaint. And no learning occurs.

Our three contributions are grounded directly in this failure. First, we provide the empirical evidence: SST-2 accuracy of 0.5092 across three identical consecutive epochs on Mamba-130m, with MNLI degrading from 0.3463 to 0.3234 --- both below the MUST\_WORK gate and below any threshold for meaningful learning. Second, we document verified zero-shot GLUE baselines for \texttt{mamba-130m-hf} that are independently reusable by any future Mamba adaptation experiment. Third, we characterize the diagnostic signature --- loss oscillation without monotonic decrease, combined with flat per-epoch accuracy --- so practitioners can detect this failure mode during training.

Three directions follow directly from the failure analysis. Exact replication of the MambaPEFT result (90--92\% SST-2) is the highest-priority step, following the ranked ablation order (checkpoint type first). If adaptation is needed at the scan level, \texttt{dt\_proj} LoRA and direct $B/C$ parameter adaptation operate inside or adjacent to the SSM computation --- the theoretically motivated next step. And prefix tuning bypasses the scan entirely by adapting at the input embedding level.

The SSM case is not a dead end. Standard LoRA asks Mamba to learn through a wall. The contribution of this work is measuring the wall, describing its signature, and pointing to the door --- architecture-aware PEFT methods that route through the SSM computation rather than around it.
